\section{总结与延伸}

\subsection{核心观点总览}
\begin{table}[H]
\centering
\begin{tabular}{p{2.8cm}p{4.0cm}p{5.0cm}}
\toprule
\textbf{主题} & \textbf{panel 结论} & \textbf{对应系统问题} \\
\midrule
推理 Infra & PD 分离、大 EP 已经进入生产主流 & 如何按 workload 重新组织资源？ \\
RL Infra & rollout 是核心瓶颈，训推结合越来越重要 & 如何让大模型 RL 可扩展？ \\
Agent Infra & 环境、workflow、事件流使系统更动态 & 如何承接真实工具与真实失败？ \\
训练系统 & 从单主线框架走向异构调度系统 & 如何模块化拆出 infra 问题？ \\
未来方向 & 不只是训练模型，而是在训练系统 & 如何让系统和算法共同迭代？ \\
\bottomrule
\end{tabular}
\caption{Infra 专题最值得记住的五条线索}
\end{table}

\subsection{三条实践导向的 takeaway}
\begin{enumerate}
  \item 如果工作负载已经变成 Agent / RL / multi-modal，就不要再用纯 pretraining 的眼光设计系统。
  \item 真正限制大模型能力上限的，越来越多是 rollout、environment、调度和故障恢复，而不是单一算子速度。
  \item 基础设施成熟的标志，不是用户感知到它很复杂，而是它能在复杂 workload 下保持稳定、可复现、可扩展。
\end{enumerate}

\begin{importantbox}{Infra 专题的真正启发}
它提醒我们，大模型时代的 ``系统'' 不是附属品，而是能力的一部分。很多看起来是算法瓶颈的问题，只有在系统设计、环境构造和部署接口都跟上时，才可能真正被解决。
\end{importantbox}

\subsection{拓展阅读}
\begin{itemize}
  \item MoonCake、KTransformers、vLLM、SGLang 等推理系统资料
  \item RL infra 与 rollout engine 相关开源项目和论文
  \item Agent environment / world model / benchmark 构造方向工作
  \item 分布式训练、容错恢复、文件系统与数据库在 AI Infra 中的迁移经验
  \item 关于 heterogeneous workflow 和 event-driven AI systems 的工程实践
\end{itemize}

\end{document}
